\documentclass[11pt,a4paper]{article}
\usepackage[T1]{fontenc}
\usepackage[utf8]{inputenc}
\usepackage{mathpazo}
\usepackage[scaled=0.92]{helvet}

\usepackage{microtype}
\usepackage{graphicx}
\usepackage[font=small,labelfont=bf]{caption}
\renewcommand{\thefigure}{20.\arabic{figure}}
\usepackage[margin=2.6cm]{geometry}
\usepackage[round]{natbib}
\usepackage{titlesec}
\usepackage{enumitem}
\usepackage{xcolor}
\usepackage[most]{tcolorbox}
\usepackage[hidelinks]{hyperref}
\usepackage{url}
\urlstyle{same}

\setlist[itemize]{itemsep=2pt,topsep=4pt}
\setlength{\parskip}{6pt}
\setlength{\parindent}{0pt}
\renewcommand{\thesection}{20.\arabic{section}}

% Footnote citation in the style of the first edition
\newcommand{\fn}[1]{\footnote{\citealp{#1}.}}

% Claim-status marks (underlines, margin bars and the key)
\usepackage{gpt-status}
\usepackage{longtable}
\usepackage{array}
\usepackage{booktabs}
\usepackage{float}
\usepackage{placeins}
% Explanatory registers
\newcommand{\reg}[1]{\textsf{\small\textbf{[#1]}}}
\newcommand{\RN}{\reg{N}}
\newcommand{\RGR}{\reg{GR}}
\newcommand{\RFR}{\reg{FR}}
\newcommand{\RC}{\reg{C}}
\newcommand{\REN}{\reg{EN}}
\newcommand{\RQ}{\reg{Q}}
\newcommand{\RDS}{\reg{DS}}
\newcommand{\expl}[1]{\par\noindent\hspace*{1.2em}\begin{minipage}[t]{\dimexpr\linewidth-1.2em}\small\textit{Explained as:} #1\end{minipage}\par\smallskip}
\newcommand{\why}[1]{\par\textit{#1}\par}
\usepackage{tocloft}
\renewcommand{\cftsecfont}{\bfseries}\renewcommand{\cftsecpagefont}{\bfseries}\setlength{\cftsecnumwidth}{3.2em}
\begin{document}
\renewcommand{\thetable}{20.\arabic{table}}
\begin{titlepage}
\centering
\vspace*{1.5cm}
{\Huge\bfseries Graviton Pressure Theory\par}
\vspace{0.4cm}
{\Large The Unified Framework\par}
{\large Individual Submission\par}
\vspace{0.3cm}
{\normalsize This document is part of a multi-part scientific framework\par}
\vspace{1.6cm}
{\large Part 20 of 30\par}
\vspace{0.6cm}
{\huge\bfseries Coherence\par}
\vspace{0.5cm}
{\Large The Order and Structure of Matter\par}
\vspace{1.4cm}
\begin{minipage}{0.82\textwidth}\centering\small
This submission is part of the broader Graviton Pressure Theory (GPT) project, a redefinition of gravitational interaction rooted in causal field dynamics and physical force transmission. Each document is designed to stand on its own; its full context emerges as part of the larger framework. For the complete series, see \url{https://peopleoftheworld.net/graviton-pressure-theory/}.
\end{minipage}
\vfill
{\large Author: Shareef Ali Rashada\par}
{\normalsize shareef@peopleoftheworld.net\par}
\vspace{0.3cm}
{\normalsize Second edition, November 2026\par}
\end{titlepage}

\tableofcontents
\newpage

\section*{About This Paper}
\addcontentsline{toc}{section}{About This Paper}

Part~19 said what the states of matter are: degrees of collective phase order among the knots of a body, from the fixed arrangement of a crystal to the shared phase of a superconductor.\fn{rashadaPart19} This paper is about the order and structure of matter, and above all how matter changes its order: how each level of order is held, what heat does to it, why some substances lose their order at one temperature and others over a range, with a definite exchange of energy, why it forms again on cooling and what it keeps when it does, and what each kind of order allows, from the stiffness of a solid to the flow of a liquid with no friction at all.

Its central idea is the first edition's, and it stands stronger now than when it was written: \argm{heat does not destroy structure by force; it destroys the timing that order requires.} Part~18 supplied the mechanism, rhythms locking through what they share, and Part~19 supplied what is locked, the phases of the knots.\fn{rashadaPart18} The graviton is in Part~15, the field in Part~16 and mass in Part~17.\fn{rashadaPart15}\fn{rashadaPart16}\fn{rashadaPart17} Claims are marked by their status, as the key shows.

\statuskey

\section{Order Has Levels}
\label{sec:levels}

\begin{argbar}
Matter was built in levels as the universe cooled (Part~19): knots, then nucleons and nuclei, then atoms, then bonds between atoms, then arrangements of many atoms. \textbf{Heating takes the levels off again in the reverse order, each at its own measured cost.} What a body loses at a transition is one level of order; everything beneath it is kept.
\end{argbar}

Iron shows the whole ladder (Figure~\ref{fig:ladder}). At 1,043 kelvin its aligned spins disorder and it stops being a magnet. At 1,811 kelvin it melts, taking in 13.8 kilojoules a mole, 0.14 electronvolts an atom. At 3,134 kelvin it boils, taking in about 340 kilojoules a mole, 3.5 electronvolts an atom, twenty-five times as much. Hotter again, near ten thousand kelvin and beyond, its atoms lose their outer electrons, 7.9 electronvolts for the first. Its nuclei hold together until energies of 8.8 million electronvolts a nucleon, and the nucleons themselves dissolve into a single liquid of their parts only near two trillion kelvin, in the collisions of heavy nuclei at the largest accelerators.\fn{bazavov2019} The knots at the bottom of it all, the electron and the proton's lobes, are not undone by heat at all; a knot is unmade only by meeting its mirror (Part~19).

\begin{figure}[tb]
\centering
\includegraphics[width=\textwidth]{fig20-1-ladder.pdf}
\caption{The ladder of order, read upward: the energy per particle needed to undo each level, with the same energy as a temperature on the top scale. Each step costs more than the one beneath it, by factors from about two to a million. Heating removes the levels from the bottom of the figure upward; the cooling universe built them in the opposite order (Part~19).}
\label{fig:ladder}
\end{figure}

\argm{The ladder answers a question the first edition raised.} It said that ``melting is not caused by bond dissociation, but by the structural collapse of graviton coherence''. The numbers agree. The energy of iron's bonds, its cohesive energy, is 4.3 electronvolts an atom, about $5.8\times10^{10}$ joules in every cubic metre; melting takes 0.14 electronvolts an atom, about 3\% of it. So melting removes the long-range arrangement of the atoms and leaves their bonds almost whole; it is boiling that breaks the bonds. A liquid is a solid that has lost its arrangement and kept its bonds.

The same reading corrects a gas. The first edition called a gas ``fully decoherent'', with ``all phase-aligned oscillation'' lost. A gas has lost only the order \emph{between} its atoms; every atom in it keeps all of its own internal order, its electron patterns locked around its nucleus, its nucleons bound, its knots intact. Coherence is nested, and a transition is the loss of one level of it. \argm{Some substances show the levels coming off one at a time.} A liquid crystal such as 8CB is a crystal below 21.5\,$^\circ$C; above it, its rod-shaped molecules lose their fixed positions but stay in layers; above 33.5\,$^\circ$C they lose the layers but still point the same way; above 40.5\,$^\circ$C they lose that too and the substance is an ordinary liquid. Three transitions, each removing one kind of order and keeping the rest.

\section{What Heat Is}
\label{sec:heat}

\argm{Heat is energy delivered as disordered motion.} It reaches a body in two ways, and each can be followed to what touches what.
\begin{itemize}
\item \textbf{By knocks.} At a contact the patterns of two atoms overlap, and patterns of the same circulation push each other apart (exclusion, Part~19). A fast atom striking a slow one hands over momentum through that push. Each knock shifts the struck pattern's motion and its phase relative to its neighbours.
\item \textbf{By light.} Light is a transverse wave of the medium (Part~15). Its energy is taken up by the electron patterns, which shift into higher standing patterns and pass the energy to their neighbours as they settle, again as knocks.
\end{itemize}

\hypm{So heat acts on order through phase.} The first edition said that heat is ``not an energizer but a decoherence catalyst''. Measurement requires both readings at once. Heat is energy: ice takes in 334 kilojoules a kilogram to melt, and a body's mass rises by the energy it holds divided by $c^2$, so a kilogram of iron melted is heavier by $2.8\times10^{-12}$ of its mass (Part~17). And what that energy does is disorder: it scatters the phases that order keeps aligned. \argm{The energy is how much is delivered; decoherence is what it does.} Temperature measures the average energy of the disordered motion, about $kT$ for each way a particle can move, where $k$ is Boltzmann's constant. It is a measure of the knocking, not a separate cause.

\section{Locking against Knocks}
\label{sec:locking}

\begin{hypbar}
Part~18 showed that coupled rhythms lock: each nudges the other through something they share, and if the nudging is stronger than the difference between their rates, they keep step. In GPT the coupling between the patterns of a body is carried by the medium, through their overlap and exchange. Heat adds random knocks to every rhythm. \textbf{Order holds while locking restores step faster than knocks destroy it; a transition is where the knocks win.}
\end{hypbar}

This is a known piece of mathematics before it is a claim about matter. In 1975 Yoshiki Kuramoto described a population of rhythms, each with its own natural rate, each pulled toward the common phase of all.\fn{kuramoto1975} Below a critical strength of coupling the rhythms run independently; above it a common phase appears suddenly and grows. Adding random knocks has the same effect as weakening the coupling: with the coupling fixed, order falls as the knocks strengthen and vanishes at a definite point (Figure~\ref{fig:locking}). The same mathematics describes flashing fireflies, the pacemaker cells of the heart and arrays of electrical oscillators; Part~18 named each as an example of locking.

\begin{figure}[tb]
\centering
\includegraphics[width=\textwidth]{fig20-2-locking.pdf}
\caption{A model of locking against knocks: 500 coupled rhythms with a spread of natural rates and random knocks added. Left: the order, how far their phases keep step, as the knocks strengthen; it falls away near the point where the knocks outpace the coupling, as the model's theory predicts. Right: each oscillator's phase as a dot on a circle, locked under weak knocks and spread round the circle under strong ones.}
\label{fig:locking}
\end{figure}

The first edition already had the measure of this order. Its part on magnetism (Part~22 in this edition) defined the coherence of $n$ spins with phases $\phi_i$ as
\begin{equation}
C_s = \frac{1}{n}\left|\sum_{i=1}^{n} e^{i\phi_i}\right|,
\end{equation}
which is 1 when every phase is the same and near 0 when they are scattered round the circle. It is the same quantity Kuramoto's model uses to measure how far a population keeps step, and it is what Figure~\ref{fig:locking} plots.

\argm{Matter shows the signs of this contest in measured rules.}
\begin{itemize}
\item \textbf{Lindemann's rule (1910).} A solid melts when the jiggling of its atoms reaches about a tenth of the distance between them.\fn{lindemann1910} The first edition's version: collapse comes ``when internal oscillations can no longer align''. It was already measured.
\item \textbf{Richards's rule (1897).} For metals, the energy taken in on melting is close to $kT$ at the melting point for every atom: iron takes 0.143 electronvolts an atom at 1,811 kelvin, where $kT$ is 0.156.\fn{richards1897} Melting costs about one typical knock's worth of energy per atom. \hypm{In GPT that is the price of releasing each atom from the common phase of the lattice.}
\item \textbf{Trouton's rule.} Boiling takes about ten times $kT$ per molecule at the boiling point for most liquids, and somewhat more for metals (iron, thirteen). Bonds cost an order of magnitude more to undo than arrangement.
\end{itemize}

\paragraph{How order is lost depends on how it is held.} Some substances lose their order at one temperature, and others over a range. Which one happens is set by the structure, not by the kind of transition.
\begin{itemize}
\item \textbf{At one temperature:} pure, uniform crystals such as metals, ice and salt. At its melting point a solid takes in its latent heat at a fixed temperature, and while it does, solid and liquid sit side by side.
\item \textbf{Over a range:} glasses, waxes, butter and many plastics have no melting point; they soften gradually. Most alloys melt between two temperatures: common 60/40 tin--lead solder is pasty between 183 and 190\,$^\circ$C, while the one exact mixture called the eutectic, 63/37, melts sharply at 183.
\item \textbf{Different arrangements of one substance:} cocoa butter can crystallise in six forms, which melt at six temperatures between about 17 and 36\,$^\circ$C. Same molecules, different order, different melting point.
\item \textbf{Fewer neighbours, weaker hold:} even a sharply melting crystal begins to melt at its surfaces and grain boundaries, where atoms have fewer neighbours. Ice carries a liquid-like layer at its surface below 0\,$^\circ$C, as Faraday proposed in 1850, and gold particles a few nanometres across melt hundreds of degrees below bulk gold.\fn{buffat1976}
\end{itemize}

\hypm{In GPT the reason is how each pattern is held in step: by its neighbours, through the medium.} In the mathematics of locking, the pull that brings each rhythm back into step grows with how many of the others are already in step. So the lock reinforces itself, and its loss can cascade: when some patterns let go, their neighbours lose part of their hold and follow. \hypm{Where every pattern is held the same way, as in a uniform crystal, all reach their limit together and the loss cascades through the whole body at one temperature. Where the hold varies from place to place, as in a glass, a mixture or a body with many surfaces, each region lets go at its own temperature and the loss spreads over a range.} The small particles show that the hold depends on the neighbours: fewer neighbours, a weaker hold, a lower melting point. How the hold is carried inside a body, what flows through its structure and what directs that flow, is the subject of Part~21.\fn{rashadaPart21}

\section{Cooling: Order Forms, and Needs a Seed}
\label{sec:cooling}

When a body cools through a transition, the order forms again and the energy comes back out: freezing water gives up the same 334 kilojoules a kilogram it took to melt it. \hypm{In GPT the released energy is the disorder handed back to the medium, as light and as knocks to the surroundings, when the patterns fall into step.}

\argm{Order needs somewhere to start.} Clean water can be cooled far below its freezing point and stay liquid; small droplets of pure water reach about $-38$\,$^\circ$C before they freeze on their own.\fn{koop2000} A dust grain, a scratch or a crystal of ice dropped in starts the freezing at once, and it spreads from that point. In the language of locking, the common phase needs a first group of patterns already in step, large enough that it pulls its neighbours in faster than knocks pull them out. \hypm{A seed is a region already locked, and freezing is the lock spreading.}

\argm{Order formed on cooling keeps a record of the conditions in which it formed.} Lava cooling through iron's Curie point is magnetised in the direction of the Earth's field at that moment and keeps it for hundreds of millions of years; the record of the field's reversals in the ocean floor is read this way.\fn{vine1963} A magnet's magnetisation depends on its history as well as on the field it is in now, the effect called hysteresis. Part~18 said that the past exists as persisting structure. \argm{Memory in matter is literal: it is order that was set at a transition and has not since been undone.} The first edition's phrase ``heat erases memory'' is kept in this sense. Part~13's activation record proposes the same principle for living things: a small bias at a moment of ordering decides a lasting order.\fn{rashadaPart13}

\section{Solids: Arrangement Held in Step}
\label{sec:solids}

\begin{hypbar}
A solid is a body whose atoms keep a fixed, long-range arrangement. In GPT each atom is a standing pattern of the medium, locked around its nucleus (Part~19). In a solid these patterns are also locked to one another:
\begin{itemize}
\item \textbf{Bonds} are patterns overlapping in spatial phase, reinforcing between the atoms, shared by two electrons of opposite circulation (Part~19).
\item \textbf{Stiffness and hardness} come from exclusion: patterns of the same circulation cannot occupy one place, and resist being pushed together (Part~15).
\item \textbf{The arrangement} is the common phase of the whole lattice, renewed every turn of every pattern.
\end{itemize}
\end{hypbar}

The first edition described solids as ``graviton pressure traps'', with ``closed-loop corridors'' resolving the pressure of the medium inside them. The restatement names what touches what: patterns overlapping, sharing electrons, pushing back where their circulations match. Its point survives: \argm{solidity is not the absence of motion but the active maintenance of order.} Every pattern in a solid turns at its own rhythm, $10^{20}$ times a second for an electron (Part~18), and holds its place by staying in step.

\paragraph{Two pressures.} The first edition gave the internal pressure of a solid as $10^{10}$ to $10^{12}$ pascals. That is a real quantity, but not the medium's pressure in empty space. It is the energy density of matter's binding: iron's cohesive energy is $5.8\times10^{10}$ joules a cubic metre, and its resistance to compression, its bulk modulus, is $1.7\times10^{11}$ pascals. \hypm{Read in GPT, it is the energy of the medium organised into a solid's bonds.} The medium's pressure in free space is about $4.5\times10^{-10}$ pascals, the value that passes the heat budget (Part~15). The two are different quantities and both stand. What does not stand is the first edition's statement that melting iron is a fall in internal pressure from $10^{12}$ to $10^{10}$ pascals: melting takes about 3\% of the binding energy, not 99\% (Section~\ref{sec:levels}).

\paragraph{Directions.} A crystal is not the same in every direction, and its order can fail unevenly. Many crystals expand by different amounts along different axes when heated; magnetic crystals have easy axes along which their spins prefer to point. The first edition's ``anisotropic collapse'' is kept with these as its examples.

\section{Liquids, Glass and Flow}
\label{sec:liquids}

\argm{A liquid keeps order between neighbours only, and that order is constantly broken and re-formed.} X-ray and neutron scattering show it directly: in a liquid each atom has a shell of nearest neighbours at a well-defined distance, a fainter second shell, and almost no order beyond.

\begin{hypbar}
\textbf{Viscosity is the time order takes to re-form.} When a liquid is sheared, each atom's cage of neighbours is distorted; it resists as a solid would until the cage breaks and re-forms in a new place. The resistance is the stiffness of the local order multiplied by the time it takes to re-form, Maxwell's relation $\eta = G\tau$. This is Part~18's relaxation, the rate at which order returns, applied to matter.
\end{hypbar}

The relation runs across an enormous range. In water the local order re-forms in about a picosecond, and its viscosity is a thousandth of a pascal-second. Pitch looks solid and shatters under a hammer, but it flows: in the experiment begun at the University of Queensland in 1927, a drop falls from a funnel about once a decade, and its viscosity is about $2\times10^{8}$ pascal-seconds.\fn{edgeworth1984} \argm{A glass is a liquid that ran out of time to arrange itself.} As a liquid that does not crystallise cools, its local order takes longer and longer to re-form, until at about $10^{12}$ pascal-seconds the re-forming takes minutes and the liquid is held in place, disordered but fixed. That is the glass transition. The first edition's three ``exotic'' cases each now have a mechanism:
\begin{itemize}
\item \textbf{Glass:} order held without arrangement, because the time to arrange grew longer than the cooling allowed.
\item \textbf{Viscoelasticity:} a material that acts as a solid for pushes faster than its relaxation time and as a liquid for slower ones.
\item \textbf{Superfluidity:} flow with no viscosity at all, because the whole body shares one phase and there are no neighbours to slip past (Section~\ref{sec:shared}).
\end{itemize}

\section{Gases and Plasmas}
\label{sec:gases}

\argm{In a gas the atoms or molecules keep no lasting relation between their phases}; each moves freely between collisions and is knocked into a new motion at each one. What it keeps is everything inside it (Section~\ref{sec:levels}).

\argm{A plasma is a gas in which the outer electron patterns have come loose from their nuclei.} The first edition called it ``hyper-decoherence''. That is half right. One level of order is lost, the locking of electron patterns around nuclei, at the measured ionisation energies: 13.6 electronvolts for hydrogen, 7.9 for iron's first electron. But loose charges organise themselves strongly through what they carry, and plasmas are full of collective order: waves, filaments, sheets, and the reconnection of magnetic fields that drives solar flares. \hypm{A plasma loses one level of order and gains a new collective order through charge}, the hand of the knot (Part~19).

\section{Aligned Spins and the Curie Point}
\label{sec:curie}

\argm{In a magnet the spins of unpaired electrons point together across regions called domains.} Above the Curie point, 1,043 kelvin for iron, 627 for nickel and 1,388 for cobalt, the knocks outpace the coupling that aligns them and the alignment is lost (Section~\ref{sec:locking}). Below it, the alignment returns as the metal cools.

\argm{Spin is real rotation, exchanged with the body that holds it.} In 1915 Einstein and de Haas magnetised a suspended iron rod and it turned: aligning its electrons' spins took angular momentum that the rod gave up as rotation. Barnett showed the reverse the same year: spinning a rod of iron magnetises it.\fn{einstein1915} \hypm{In GPT this is expected: spin is the circulation of a knot (Part~19), and aligning spins aligns circulations.}

\begin{hypbar}
\textbf{A magnet has two locks.} One is the arrangement of its atoms, the lattice; the other is the alignment of its electrons' circulations, the spins. They are separate levels, which is why iron loses its magnetism at 1,043 kelvin, long before it melts at 1,811. In iron the alignment fades continuously to nothing at the Curie point, with no latent heat, because the lattice holds while the spins let go. \textbf{Where the two locks depend on each other, the loss is abrupt}: in manganese arsenide near 40\,$^\circ$C, and in iron--rhodium near 80\,$^\circ$C, the magnetic order and the arrangement of the lattice change together, at one temperature, with latent heat.\footnotemark
\end{hypbar}
\footnotetext{\citealp{gutfleisch2016}.}

Part~19 proposed what the alignment does: \hypm{each spinning knot's counterflow carries a swirl, and swirls aligned across a material add, so that magnetism is the collective swirl of aligned matter.} The first edition described magnetism as the ``externalization'' of graviton pressure, with the strength of a magnet growing as the cube of a ``coherent graviton density''. The cube law had no derivation and is withdrawn, and so is the statement that magnetism is gravity resolved along an axis: torsion balances carrying masses with $10^{23}$ aligned electron spins show no change in their gravity.\fn{heckel2008} Gravity and magnetism are different modes of the same medium, inflow and swirl. The first edition's magnetism parts also gave spin alignment a role in directing the medium's path through a material, and named the timing that requires. That account, with the full treatment of magnetism, belongs to Parts~21 and 22.\fn{rashadaPart22}

\section{Shared Phase}
\label{sec:shared}

\begin{argbar}
The coldest order is the most complete: many particles moving in \emph{one} phase across a whole body. It appears in three measured forms.
\begin{itemize}
\item \textbf{Superconductors} carry current with no resistance at all, below 9.3 kelvin for niobium and about 93 for some copper oxides. A current set going in a superconducting ring has been watched for years without measurable decay.
\item \textbf{Superfluid helium}, below 2.17 kelvin, flows through channels too narrow for any ordinary liquid, with no viscosity, and climbs the walls of its container as a film.\fn{kapitza1938}\fn{allen1938}
\item \textbf{Atomic condensates}, clouds of atoms cooled to billionths of a kelvin, fall into a single shared state (Part~19).
\end{itemize}
\end{argbar}

\argm{In GPT phase is literal}, the medium's own phase in each pattern (Part~19), and shared phase is the medium in a body keeping one step throughout it. Measurement shows that phase directly. In a superconducting ring the magnetic flux trapped inside comes only in whole units of $h/2e$, Planck's constant over twice the electron's charge.\fn{deaver1961} The standard account explains this as the shared phase having to close on itself around the ring a whole number of times. \hypm{In GPT this is literal: the medium's phase around the ring must return to itself, so the swirl it encloses comes in whole turns.} The $2e$ shows that the current is carried by electrons in pairs. \hypm{In GPT a pair is two knots of opposite circulation sharing one pattern, as two electrons share a bond} (Part~19); the pairs in ordinary superconductors are indeed made of electrons of opposite spin.

\paragraph{Superconductors do not reflect gravity.} The first edition said that superconductors reflect graviton fields, and read the Meissner effect, the expulsion of a magnetic field from a superconductor, as ``corridor reflection''. Superconductors weigh normally: superconducting gravimeters float a niobium sphere in a magnetic field and measure the Earth's gravity through it to about one part in $10^{11}$.\fn{goodkind1999} The claim is withdrawn. The Meissner effect is measured as currents circulating within a thin layer at the surface, about 40 nanometres deep in niobium, which cancel the field inside.\fn{meissner1933} \hypm{In GPT a magnetic field is a swirl of the medium (Part~16), and a body in one shared phase cannot hold a swirl inside it without breaking that phase; its surface circulates against the swirl and keeps it out.}

\section{Universality}
\label{sec:universality}

\argm{The first edition's central claim, that the loss of magnetism and the behaviour of fluids have one cause, has measured support.} Near the critical point of a liquid, where liquid and gas become indistinguishable, the difference between their densities shrinks as the temperature rises, following a power of the distance from the critical temperature. In 1945 Edward Guggenheim plotted the measurements for eight different substances, neon, argon, krypton, xenon, nitrogen, oxygen, carbon monoxide and methane, and they fell on one curve, with a power close to one third.\fn{guggenheim1945} Magnets whose spins can point only one of two ways lose their magnetisation near the Curie point with the same power, about 0.33.\fn{pelissetto2002} Substances that share nothing chemically, and transitions that look nothing alike, behave identically there.

The standard account explains this. Near a critical point, regions of order grow so large that the details of what is ordered stop mattering; only the dimension of space and the number of ways the order can point remain. \hypm{GPT adds what the order is: in every case a phase relation among patterns of one medium.} The first edition's instinct was to treat melting, magnetism and fluid behaviour as one subject. Universality shows that, near their critical points, they are.

\section{What Order Does Not Do}
\label{sec:notdo}

\argm{A change of order changes a body's weight only by the energy it takes in or gives out.} The total delivery of the medium to a body tracks its energy, whatever its structure (Part~15). Freezing a kilogram of water lowers its mass by $3.7\times10^{-12}$ of itself; melting a kilogram of iron raises it by $2.8\times10^{-12}$. Nothing else in the total is predicted to change. \hypm{What may change is where inside the body the delivery lands}: an ordered body may receive it differently along its planes than a disordered one, so that its centre of gravity shifts slightly from its centre of mass. Whether it does is untested, and the tests are below.

\section{Tests}
\label{sec:tests}

The first edition proposed heating a kilogram of iron through its melting point with a gravimeter 1 to 2 centimetres away, and predicted a change in gravity of ``5--10 $\mu$g''; the units it meant were microgals, $10^{-8}$ metres per second squared. The test needs reshaping. The iron's own gravity at 2 centimetres is about $1.7\times10^{-7}$ metres per second squared, and iron changes its density by a few per cent on melting, so the change of shape alone alters the reading by $10^{-9}$ to $10^{-8}$, the size of the predicted effect. The iron also passes its Curie point first, at 770\,$^\circ$C, where magnetic forces would swamp everything. And the agreed account predicts no change in the total beyond the mass change, $2.8\times10^{-12}$. Better versions:
\begin{itemize}
\item \textbf{Weigh gallium through its melting point,} at 29.8\,$^\circ$C: room temperature, no magnetism, and a mass change of $9\times10^{-13}$, far below what mass comparators resolve (about $10^{-10}$ on a kilogram). \hypm{The prediction is null in the total; any change in the weighing beyond the change of shape would be a signal.} Water through freezing is a second case.
\item \textbf{Look for a shift in where the delivery lands.} A torsion balance carrying a single crystal, turned through its orientations and taken through a change of order, looks for the centre of gravity moving away from the centre of mass. This test is shared with Part~29.\fn{rashadaPart29}
\item \textbf{Weigh a superconductor through its transition.} An experiment of this kind, named Archimedes, is being built to weigh the energy of the vacuum in a layered superconductor as it changes state.\fn{calloni2016} Its result will bound any change in weight with the onset of shared phase.
\item \textbf{One clock in two states of order.} Every kind of clock tested is slowed alike by the medium (Part~18). Whether a single clock taken through a change of its own order is slowed by the same factor is untested.
\end{itemize}

\section{What Remains to Be Derived}
\label{sec:owed}

\begin{itemize}
\item \openm{The temperatures of transitions from the phase relations of the knots: why iron melts at 1,811 kelvin and loses its magnetism at 1,043.}
\item \openm{The strength of the coupling between neighbouring patterns, carried by the medium, in terms of the medium's own properties.}
\item \openm{Why pairs form in superconductors, and the temperature at which they do, in terms of knots of opposite circulation sharing a pattern.}
\item \openm{The relaxation times of liquids, and so their viscosities, from the rate at which the medium restores order (Part~18).}
\item \openm{Whether a change of order shifts the centre of gravity of a body, and by how much.}
\end{itemize}

\section*{Conclusion}
\addcontentsline{toc}{section}{Conclusion}

Order comes in levels, built as the universe cooled and taken off again by heat in the reverse order, each at its measured cost: the arrangement of atoms first, then their bonds, then their electron patterns, then their nuclei. Heat is energy, and what it does is disorder: it knocks the patterns of a body out of step. Order holds while locking through the medium restores step faster than the knocks destroy it, and a transition is where the knocks win. How the order is held decides how it goes: all at once where every pattern is held alike and the loss cascades, over a range where the hold varies. Cooling reverses it, given a seed, and the order it forms keeps a record of the moment it formed. A solid holds its arrangement in step; a liquid keeps its bonds and loses its arrangement, and its viscosity is the time the arrangement takes to re-form; a glass ran out of that time; a gas keeps everything but the relations between its atoms; a plasma loses its electrons' locking and gains order through charge. The coldest order is one phase shared across a whole body, and in a superconducting ring that phase can be counted. Near their critical points, transitions as different as boiling and the loss of magnetism behave identically, because in every case what changes is a phase relation. Heat does not destroy; it erases memory, and memory in matter is order that has not yet been undone.

\section*{Changes in the Second Edition}
\addcontentsline{toc}{section}{Changes in the Second Edition}

This list records what changed from the first edition and why. Changes to the argument alter what the paper claims or how it supports a claim; clarifications improve wording, sources or presentation and leave the argument as it was.

\subsection*{Changes to the argument}
\begin{itemize}
\item The central idea is kept and given its mechanism: heat destroys the timing order requires; order holds while locking through the medium beats the knocks (Part~18), with Kuramoto's model, Lindemann's rule and Richards's rule as its measured anchors.
\item New: order comes in levels, and heating removes them in the reverse order of the cooling that built them (Part~19), each at its measured cost. Melting takes about 3\% of iron's bond energy, which supports the first edition's statement that melting is not the breaking of bonds.
\item ``Heat is not an energizer'' becomes ``heat is energy, and what it does is disorder'': latent heat and the change of mass by $\Delta E/c^2$ are measured.
\item ``Gases are fully decoherent'' becomes: a gas loses only the order between its atoms. Plasma is restated as one level lost and a new collective order through charge gained.
\item Solidity is restated in terms of bonds, exclusion and common phase (Part~19), in place of pressure traps and closed-loop corridors. The internal pressure of $10^{10}$--$10^{12}$ pascals is identified as the energy density of binding, distinct from the medium's free-space pressure; the fall from $10^{12}$ to $10^{10}$ pascals at melting is withdrawn.
\item New: viscosity as the time order takes to re-form (Maxwell's relation), with glass, viscoelasticity and superfluidity each given a mechanism.
\item New: cooling, seeds and memory: supercooling, nucleation, the magnetic record in lava, hysteresis.
\item New: how order is lost depends on how it is held. Uniform crystals lose it at one temperature, because the lock reinforces itself and its loss cascades; glasses, mixtures, polymorphs, surfaces and small particles lose it over a range or at other temperatures. Liquid crystals show the levels coming off one at a time.
\item New: a magnet has two locks, the lattice and the alignment of circulations; where they depend on each other the loss is abrupt. Spin as real rotation (Einstein and de Haas, Barnett). The first edition's coherence measure $C_s$ is identified as the order parameter of locking.
\item Magnetism as the externalization of graviton pressure, the cube law $P_g\propto n_g^3$ and ``magnetism is gravity resolved anisotropically'' are withdrawn; magnetism is restated as the collective swirl of aligned matter and moved to Part~22.
\item ``Superconductors reflect graviton fields'' is withdrawn: superconductors weigh normally. The Meissner effect is restated as a hypothesis about a body in shared phase keeping out a swirl. New: flux quantisation as the medium's phase closing on itself, and pairs as knots of opposite circulation.
\item New: universality as measured support for the first edition's claim that these transitions share one cause.
\item The gravimetric test at the melting of iron is reshaped: the predicted size was comparable to the change of shape alone, and the iron passes its Curie point first. It is replaced by weighing gallium, a torsion test for a shift in where the delivery lands, the Archimedes superconductor weighing and a clock in two states of order.
\end{itemize}

\subsection*{Clarifications}
\begin{itemize}
\item The title changes from \emph{Graviton Coherence} to \emph{Coherence}, and the subtitle from \emph{Field-Based Phase Transitions and the Collapse of Coherence} to \emph{The Order and Structure of Matter}. This part was Part~19 in the first edition.
\item The unit of the predicted signal is corrected from micro-g to microgal.
\item ``Corridors'' and ``graviton pressure fields'' are replaced by the terms of Parts 15 to 19; the cross-references are renumbered.
\item The lines on purpose in matter (``alignment becomes intentional'', ``matter speaks'') are moved to the second register, beside Part~2.
\item Two figures, status marks and this list are added.
\end{itemize}

\newpage
\bibliographystyle{plainnat}
\bibliography{part20}

\end{document}
